% ============================================================
% APPENDIX: Literature review details — 6 slides
% ============================================================

% ----------------------------------------------------------
% A-LR1: Full prior art comparison table
% ----------------------------------------------------------
\begin{frame}[label=app-lit-table]{Prior art comparison (full table)}

  {\scriptsize
  \begin{tabular}{l cccccc}
    \toprule
    & \makecell{$\Tavg$\\[-2pt]${\approx}\,T_0$}
    & \makecell{Local\\[-2pt]$\Sab(\rr)$}
    & \makecell{Cauchy\\[-2pt]$\Aab/4$}
    & \makecell{$\eta$\\[-2pt]$=$ AO}
    & \makecell{Compl.\\[-2pt]formula}
    & \makecell{Expos.\\[-2pt]op.\ $\QQ$} \\
    \midrule
    Kodera 2024         & $\circ$ & ---         & $\circ$     & ---   & ---     & --- \\
    Funahashi 2018      & ---     & $\triangle$ & ---         & ---   & ---     & --- \\
    Li 2019             & $\circ$ & $\triangle$ & ---         & ---   & ---     & --- \\
    Samaras \& Kuster 19 & $\triangle$ & ---     & ---         & ---   & ---     & --- \\
    Diao 2021           & $\circ$ & ---         & ---         & ---   & ---     & --- \\
    Diao 2024           & $\circ$ & ---         & ---         & ---   & ---     & --- \\
    Sasaki 2017         & $\triangle$ & ---     & ---         & ---   & ---     & --- \\
    Bamba 2012, 2014    & $\circ$ & ---         & ---         & ---   & ---     & --- \\
    Flintoft 2014       & ---     & ---         & $\circ$     & ---   & ---     & --- \\
    Zhang \& Robinson 20 & $\circ$ & ---         & $\circ$     & ---   & ---     & --- \\
    Azzam 2015           & $\bullet$ & ---       & ---         & ---   & ---     & --- \\
    Gosselin 2011       & ---     & ---         & ---         & ---   & $\circ$ & --- \\
    Hochwald 2014       & ---     & ---         & ---         & ---   & ---     & $\circ$ \\
    Ying 2015, 2017     & ---     & ---         & ---         & ---   & $\circ$ & $\bullet$ \\
    Castellanos 2020    & ---     & $\triangle$ & ---         & ---   & ---     & $\triangle$ \\
    \midrule
    \textbf{This work}   & $\bullet$ & $\bullet$ & $\bullet$   & $\bullet$ & $\bullet$ & $\bullet$ \\
    \bottomrule
  \end{tabular}
  }

  \vspace{4pt}
  {\tiny
    $\bullet$ = derived analytically \quad $\circ$ = observed empirically \quad
    $\triangle$ = partial \quad --- = not addressed
  }

  \vspace{2pt}
  \returntoindex{sec:part1}{literature slide}\quad\returntosummary{sum-lit}
\end{frame}

% ----------------------------------------------------------
% A-LR2: Framework predictions vs literature
% ----------------------------------------------------------
\begin{frame}[label=app-lit-predictions]{Framework predictions vs literature}

  {\scriptsize
  \begin{tabularx}{\textwidth}{XXl}
    \toprule
    Published observation & Framework prediction & Agr. \\
    \midrule
    Bamba 2012: $\eta \approx 0.5$ at 1.5--6\,GHz
    & $\eta = T_0$; pseudo-Brewster explains weak $f$-dep.
    & 2--4\,\% \\[3pt]
    Kodera 2024: ACS $\propto$ projected area $>$ 10\,GHz
    & $\Tavg \approx T_0 \Rightarrow P_{\mathrm{abs}} = T_0 \Aperp \Sinc$
    & $<5\,\%$ \\[3pt]
    Diao 2024: $T \approx 0.52$ at 28\,GHz (FDTD)
    & Normal-incidence Fresnel: $T_0 = 0.536$
    & 3\,\% \\[3pt]
    Flintoft 2014: $\langle Q^a\rangle \approx 0.41$, 7--11\,GHz
    & $T_0 \times \eta \approx 0.48 \times 0.85 = 0.41$
    & 3--8\,\% \\[3pt]
    Zhang 2020: $\xi \approx 0.45$--0.65 above 6\,GHz
    & $\xi = \Tbar \times \Aab/A$; predicted 0.43--0.49
    & Within scatter \\[3pt]
    Li 2019: flux insensitive to angle, 30\,GHz--1\,THz
    & Pseudo-Brewster: $\Tavg/T_0 \approx 1$ for $|\ntilde|>2.5$
    & $< 5.6\,\%$ \\[3pt]
    Diao 2021: heating factor $\approx$ const vs angle
    & TE decreases, TM increases; $\Tavg \approx T_0$
    & $<6\,\%$ \\
    \bottomrule
  \end{tabularx}
  }

  \vspace{4pt}
  \returntoindex{sec:part1}{literature slide}
\end{frame}

% ----------------------------------------------------------
% A-LR3: Dosimetry prior art
% ----------------------------------------------------------
\begin{frame}[label=app-lit-dosimetry]{Prior art: local absorption and dosimetry}

  \textbf{Key finding}: all observe $T_{\mathrm{tr}}$ numerically; none derive the mechanism.

  \vspace{6pt}
  {\small
  \begin{itemize}
    \item \textbf{Kodera 2024}: $\SARwb = T_{\mathrm{tr}} \times \Aperp \times \Sinc/W$ from 1D slab. Reproduces 3D FDTD to 5\,\%, 10--100\,GHz. No physical derivation.
    \item \textbf{Funahashi 2018}: $\Sinc \times T_{\mathrm{tr}}$ correlates with temperature rise to 15\,\%, 30--300\,GHz. Normal incidence only.
    \item \textbf{Li 2019}: transmitted flux insensitive to angle/polarisation, 6\,GHz--1\,THz. Unpolarised average $\Tavg$ not examined.
    \item \textbf{Samaras \& Kuster 2019}: TE/TM transmittances vs angle. TM worst case 0.92 at pseudo-Brewster angle (exceeds ICNIRP by 70\,\%). Unpolarised average not evaluated.
    \item \textbf{Diao 2024}: $T \approx 0.52$ at 28\,GHz from full FDTD. Fresnel: $T_0 = 0.536$ (3\,\% discrepancy).
    \item \textbf{Azzam 2015}: proved $|\ntilde| > 2.5$ $\Rightarrow$ reflectance $\approx$ const over 0--60$^\circ$. Optics; not applied to dosimetry.
  \end{itemize}
  }

  \vspace{4pt}
  \returntoindex{sec:part1}{literature slide}
\end{frame}

% ----------------------------------------------------------
% A-LR4: Diffuse-field prior art
% ----------------------------------------------------------
\begin{frame}[label=app-lit-diffuse]{Prior art: diffuse-field dosimetry}

  \textbf{Key finding}: Cauchy's projected area and a transmission coefficient appear as separate empirical ingredients. Not connected to Fresnel theory.

  \vspace{6pt}
  {\small
  \begin{itemize}
    \item \textbf{Andersen 2007}: room electromagnetics. $P_{\mathrm{abs}} = S_{\mathrm{total}} \langle\sigma_a\rangle$. 153 citing papers.
    \item \textbf{Bamba 2012, 2014}: whole-body SAR in reverberation chamber at 1.8\,GHz. $\eta \approx 0.5$ from FDTD (1.45--5.8\,GHz). No physical origin identified.
    \item \textbf{Flintoft 2014}: ACS of 60 volunteers, 1--12\,GHz. Normalised ACS $\langle Q^a\rangle \approx 0.41$ at 7--11\,GHz. Dip below 3\,GHz attributed to fat-layer standing waves; no quantitative model.
    \item \textbf{Zhang 2017, 2020}: ACS of 48 subjects, 1--18\,GHz. $\xi \approx 0.45$--0.65 above 6\,GHz. ACS--BSA correlation $> 0.9$. Physical content $\xi = \Tbar \times \Aab/A$ not derived.
    \item \textbf{Hallbj\"orner 2005}: ACS of large lossy objects depends only on permittivity, not shape, in skin-depth limit. Mechanism (pseudo-Brewster) not identified.
  \end{itemize}
  }

  \vspace{4pt}
  \returntoindex{sec:part1}{literature slide}
\end{frame}

% ----------------------------------------------------------
% A-LR5: Coherent MIMO prior art
% ----------------------------------------------------------
\begin{frame}[label=app-lit-mimo]{Prior art: coherent MIMO and exposure}

  \textbf{Key finding}: all prior work has FDTD-calibrated exposure matrices. No closed-form physical content.

  \vspace{6pt}
  {\small
  \begin{itemize}
    \item \textbf{Hochwald 2014}: SAR matrix $\mathbf{S}$ at 1.8\,GHz. SAR-aware coding: +2.5\,dB over Alamouti. $\mathbf{S}$ calibrated numerically per device.
    \item \textbf{Ying 2015}: capacity-optimal precoder $\xx^* = (\lambda\mathbf{S}+\nu\II)^{-1}\hh^*$. Strong duality proved. SU-MIMO uplink.
    \item \textbf{Ying 2017}: MU-MIMO uplink. Per-user SAR matrix. Decentralised protocol.
    \item \textbf{Castellanos 2020}: per-element Fresnel at 28\,GHz. Calibration scalar $\beta \approx 0.70$. Relation to $T_0$ not examined. No body-surface integral.
    \item \textbf{Ebadi-Shahrivar 2019}: worst-case SAR $= \lambda_{\max}(\mathbf{S})$. Compliance via eigendecomposition.
    \item \textbf{Zhou 2026}: pointwise exposure matrices from dipole radiation at 30\,GHz + thermal model. $T_{\mathrm{tr}} = 0.8$ as parameter (not derived).
  \end{itemize}
  }

  \vspace{4pt}
  {\footnotesize All assume uplink. Downlink base-station geometry not addressed.}

  \vspace{4pt}
  \returntoindex{sec:part3}{literature slide}
\end{frame}

% ----------------------------------------------------------
% A-LR6: Numerical approaches
% ----------------------------------------------------------
\begin{frame}[label=app-lit-numerical]{Prior art: numerical approaches and limits}

  FDTD remains the regulatory standard.

  \vspace{6pt}
  {\small
  \begin{itemize}
    \item Single plane wave at 28\,GHz on anatomical phantom: ${\sim}10^{12}$ mesh cells
    \item Systematic campaign: 15 frequencies $\times$ 4 phantoms $\times$ 2 scenarios = 550 simulations, weeks of GPU
  \end{itemize}
  }

  \vspace{6pt}
  \textbf{Shortcuts pursued:}

  {\small
  \begin{itemize}
    \item \textbf{Kodera 2024}: 1D multilayer slab $\rightarrow$ $T_{\mathrm{tr}}$. Reproduces 3D to 5\,\%, 10--100\,GHz.
    \item \textbf{Kushiyama \& Nagaoka 2025}: spherical-wave lookup tables (${\sim}$2.2\,TB per antenna). Field superposition without re-running FDTD.
    \item \textbf{Zhang 2017}: surface-integral-equation method (ABCD + MoM). Memory scales with surface area, not volume. Validated in 2D at 1--6\,GHz.
  \end{itemize}
  }

  \vspace{6pt}
  All retain numerical tissue-parameter input. No closed-form transmission coefficient or exposure operator.

  \vspace{4pt}
  \returntoindex{sec:part1}{literature slide}
\end{frame}
